% ============================================================
% core/c1-duality.tex —— C1 轨道五条核心词条
% 词条：电磁对偶 / Langlands 对偶群 / 拓扑扭化与典范参数 / Hitchin 模空间 / Hitchin 纤维化
% 只依赖 preamble.tex 提供的环境与宏；不引入宏包、不自定义命令、不插入外部文件、不用浮动体。
% ============================================================

% ------------------------------------------------------------
% 1. 电磁对偶
% ------------------------------------------------------------
\begin{centry}{电磁对偶}{Electric-Magnetic Duality}{em-duality}
\oneline{四维 $\mathcal N=4$ 超对称杨–米尔斯理论在强耦合与弱耦合之间的精确等价，
它把规范群换成 Langlands 对偶群、把复耦合换成其倒数——整篇 KW06 的起点。}

\begin{definition}[Montonen--Olive 对偶]
设 $G$ 为紧单连通单李群，$\mathfrak g=\mathrm{Lie}(G)$，规范理论有复耦合
\[
\tau=\frac{\theta}{2\pi}+\frac{4\pi i}{g^{2}}\in\mathbb H ,
\]
其中 $g^{2}$ 是规范耦合，$\theta$ 是瞬子数的角变量：拓扑项为
\[
I_{\theta}=-\frac{\theta}{8\pi^{2}}\int_{M}\mathrm{Tr}\,F\wedge F ,
\]
而 $(1/8\pi^{2})\int_{M}\mathrm{Tr}\,F\wedge F$ 取整数值（$M$ 为闭四维流形时），故瞬子权重是 $e^{2\pi i\tau k}=e^{i\theta k-8\pi^{2}k/g^{2}}$，且 $\theta\mapsto\theta+2\pi$（即 $\tau\mapsto\tau+1$）是精确对称。称 $(G,\tau)$ 与 $({}^{L}G,\tau')$ 电磁对偶，若两者的量子理论等价，且
\[
\tau'=-\frac{1}{n_{\mathfrak g}\,\tau},\qquad
n_{\mathfrak g}=\begin{cases}1,&G\ \text{单边（simply-laced）},\\ 2,&G=B_{n},\,C_{n},\,F_{4},\\ 3,&G=G_{2}.\end{cases}
\]
等价的内容包括规范群的整体结构，以及电、磁荷 $(e,m)$ 的线性变换。
\end{definition}

\begin{theorem}[单边情形的自对偶点]
设 $G$ 单边、$\theta=0$，则 $\tau=4\pi i/g^{2}$；$S$ 变换 $\tau\mapsto-1/(n_{\mathfrak g}\tau)=-1/\tau$ 的不动点满足 $g^{2}=4\pi$。在该点理论自对偶，且 $W$ 玻色子与磁单极子质量相等。
\end{theorem}
\begin{proof}
把 $\tau=4\pi i/g^{2}$ 代入 $S$：$-1/\tau=i g^{2}/(4\pi)=4\pi i/g'^{2}$，故 $g'^{2}=16\pi^{2}/g^{2}$，不动点方程 $g^{4}=16\pi^{2}$ 给出 $g^{2}=4\pi$。磁、电质量之比由 §2.2 的 BPS 质量公式 $M=|a\,e+a_{D}m|$ 给出（中心荷 $Z=\sqrt{2/\mathrm{Im}\,\tau}\,\vec\varphi\cdot(\vec n+\tau\vec m)$，式 (2.22)；$a$ 是电标量的期望值、$a_{D}=\tau a$，$e,m$ 为电、磁荷）：该比为 $|a_{D}/a|=|\tau|=4\pi/g^{2}$，在 $g^{2}=4\pi$ 处等于 $1$。取 $m_{W}=gv$ 的归一化，则 $m_{\text{单极子}}=m_{W}\cdot 4\pi/g^{2}=4\pi v/g$，两者在自对偶点同为 $\sqrt{4\pi}\,v$。式 (2.17) 的 $S$ 矩阵正实现 $\tau\mapsto-1/(n_{\mathfrak g}\tau)$。
\end{proof}

\physside{$S$-对偶同时作用在三样东西上：复耦合 $\tau\mapsto-1/(n_{\mathfrak g}\tau)$、规范群 $G\mapsto{}^{L}G$、线算子的电/磁荷按 $SL(2,\mathbb Z)$ 的元素线性混合。$\theta$ 项的角色有二：其一，瞬子数取整使 $\tau\mapsto\tau+1$ 为精确对称，于是对偶群 $\Gamma$ 由 $T:\tau\mapsto\tau+1$ 与 $S:\tau\mapsto-1/(n_{\mathfrak g}\tau)$ 生成；单边时 $\Gamma=SL(2,\mathbb Z)$，非单边时 $\Gamma$ 只是 $SL(2,\mathbb R)$ 的 Hecke 子群（因子 $n_{\mathfrak g}$ 来自 BPS 质量公式，反映根与余根的关系）。其二，在紧化到二维时 $\theta$ 项变成 sigma 模型的 $B$-场。KW06 §2 只用它作用在 $(\tau,G)$ 上的这一面，故把它当作「把 $G$-理论送到 ${}^{L}G$-理论」的精确对称性来用。}

\begin{example}[$G=SU(2)$，$\theta=0$]
$g^{2}=4\pi\approx12.57$，即 $g=\sqrt{4\pi}\approx3.545$。取 $v=1$：$m_{W}=gv\approx3.545$，$m_{\text{单极子}}=4\pi v/g=4\pi/\sqrt{4\pi}=\sqrt{4\pi}\approx3.545$，两者相等且都等于 $\sqrt{4\pi}\,v$。若改取 $g^{2}=4$（$g=2$），则 $m_{W}=2v$ 而 $m_{\text{单极子}}=\pi v$，并不相等——自对偶只在 $g^{2}=4\pi$ 处发生。顺带核对：$\tau=4\pi i/g^{2}=i$，确实满足 $S:\tau\mapsto-1/\tau=i=\tau$。
\end{example}

\begin{remark}[归一化与两个常见误读]
其一，KW06 把标量归一化得使动能与 $\tau$ 无关，故中心荷 (2.22) 带因子 $\sqrt{2/\mathrm{Im}\,\tau}$；换成 $m_{W}=gv$ 的物理归一化后，磁、电质量之比仍是 $|\tau|=4\pi/g^{2}$，结论不变。其二，$n_{\mathfrak g}$ 不是新增参数：它由 $G$ 的边数决定，$S$ 与 $T$ 生成的群在非单边时只是 $SL(2,\mathbb R)$ 的 Hecke 子群，故不能把 $\Gamma$ 一律写成 $SL(2,\mathbb Z)$。其三，$Sp(k)$ 与 $Spin(2k+1)$ 是仅有的一对「李代数同构而整体结构不同」的单群，故两者的 $\tau$ 可以直接比较；其他非单边群的 $\tau$ 比较需要额外的归一化约定（KW06 §2.2 p.12）。
\end{remark}

\statusline{物理陈述：Montonen--Olive 1977 提出，Sen 1994 由弦论给出强证据，至今不是数学定理。可当作可靠计算的只有：$S$ 作用在 $\tau$ 与质量谱上的代数形式（KW06 §2.1 p.9 式 (2.11) 的 $\theta$ 项、p.10 式 (2.12) 的 $\tau$、§2.2 p.11 式 (2.17) 的 $S$ 矩阵与 $n_{\mathfrak g}$、p.14--16 式 (2.18)--(2.25) 的 BPS 中心荷；Witten 2008 §2），以及它诱导的典范参数变换 $\Psi\mapsto-1/(n_{\mathfrak g}\Psi)$（KW06 §2.2 p.13，§3.5 式 (3.54)）；单边时 $\theta=0$ 的 $T:\tau\mapsto\tau+1$ 亦为精确对称。KW06 全文把对偶当作前提，本册把此条列为「词典的入口」。}

\dictionaryrow{它把规范群送到 Langlands 对偶群（\ref{core:lg}）；它作用在 $\tau$ 上又诱导扭转参数 $t$ 与典范参数 $\Psi$ 的变换（\ref{core:gl-twist}）；几何 Langlands 的物理化身见 \ref{core:gl}。}
\end{centry}

% ------------------------------------------------------------
% 2. Langlands 对偶群
% ------------------------------------------------------------
\begin{centry}{Langlands 对偶群}{Langlands Dual Group}{lg}
\oneline{把根数据里的权与余权、根与余根对换，就得到 S-对偶另一侧的规范群；这个纯组合的操作是整本词典的「换字母表」。}

\begin{definition}[对偶根数据]
设 $G$ 是 $\mathbb C$ 上连通约化群，$T\subset G$ 为极大环面。$G$ 的根数据是四元组
\[
\bigl(X^{*}(T),\,X_{*}(T),\,\Delta,\,\Delta^{\vee}\bigr),
\]
其中 $X^{*}(T)=\mathrm{Hom}(T,\mathbb C^{\times})$ 是权格、$X_{*}(T)=\mathrm{Hom}(\mathbb C^{\times},T)$ 是余权格，$\Delta\subset X^{*}(T)$ 是单根集、$\Delta^{\vee}\subset X_{*}(T)$ 是单余根集。称连通约化群 ${}^{L}G$ 为 $G$ 的 Langlands 对偶群，若其根数据为
\[
\bigl(X_{*}(T),\,X^{*}(T),\,\Delta^{\vee},\,\Delta\bigr):
\]
即 ${}^{L}G$ 的极大环面 ${}^{L}T$ 满足 $X^{*}({}^{L}T)=X_{*}(T)$，且 ${}^{L}G$ 的根恰是 $G$ 的余根、单根恰是 $G$ 的单余根。等价地，$G$ 与 ${}^{L}G$ 的权/余权格互换、根/余根互换。
\end{definition}

\begin{theorem}[对偶群的基本性质]
设 $G$ 为连通约化群，根数据如上。
(i) ${}^{L}G$ 存在，且在同构意义下唯一；(ii) ${}^{L}({}^{L}G)\cong G$；
(iii) 若 $G$ 半单，则中心与基本群对偶：
\[
Z({}^{L}G)=X_{*}(T)/\mathbb Z\Delta^{\vee}\cong\pi_{1}(G),\qquad
\pi_{1}({}^{L}G)\cong Z(G);
\]
(iv) $G$ 单连通当且仅当 ${}^{L}G$ 是伴随型（中心平凡）。
\end{theorem}
\begin{proof}
(i)(ii) 是根数据分类定理（Chevalley）的直接推论：连通约化群由其根数据在同构意义下唯一确定，而对偶四元组仍满足根数据公理，取两次对偶回到原四元组。中心的计算：$Z({}^{L}G)=X^{*}({}^{L}T)/\mathbb Z\Delta({}^{L}G)=X_{*}(T)/\mathbb Z\Delta^{\vee}$，而 $\pi_{1}(G)=X_{*}(T)/\mathbb Z\Delta^{\vee}$（$G$ 半单），故二者相等；同理 $\pi_{1}({}^{L}G)=X^{*}(T)/\mathbb Z\Delta=Z(G)$。又 $G$ 单连通 $\iff$ $X^{*}(T)=\mathbb Z\Delta$ $\iff$ ${}^{L}G$ 的权格等于其根格，即 ${}^{L}G$ 为伴随型。
\end{proof}

\physside{在物理侧，${}^{L}G$ 就是 $S$-对偶后的规范群：$(G,\tau)$ 与 $({}^{L}G,-1/(n_{\mathfrak g}\tau))$ 描述同一个量子理论，$G$ 的 Wilson 线算子对应 ${}^{L}G$ 的 't Hooft 线算子。性质 (iii) 是 §7 讨论离散通量的钥匙：四维流形 $M$ 上 $G$-丛的特征类取在 $H^{2}(M,\pi_{1}(G))$，取值 $\xi$ 的路径积分给出配分函数 $Z_{\xi}$，而 $Z({}^{L}G)\cong\pi_{1}(G)$ 意味着 $\xi$ 正是 ${}^{L}G$ 中心的「$\theta$ 角」，$Z_{\xi}$ 按 ${}^{L}G$ 的表示变换（KW06 §2.2 p.13、§7）。在几何 Langlands 中，${}^{L}G$ 的平坦丛给出 $\mathrm{Bun}_{G}$ 上的 Hecke 特征层（Frenkel 2009 §2）。}

\begin{example}[对偶群小表]
$GL_{n}\leftrightarrow GL_{n}$（自对偶）；$SL_{n}\leftrightarrow PGL_{n}$，两者维数都是 $n^{2}-1$；$Sp_{2n}\leftrightarrow SO_{2n+1}$，两者维数都是 $n(2n+1)$（$n=2$ 时 $Sp_{4}$ 与 $SO_{5}$ 都是 $10$ 维）；$E_{8}\leftrightarrow E_{8}$。再核中心：$SU(2)$ 单连通、$\pi_{1}=0$，对偶群 $SO(3)$ 中心平凡；反过来 $\pi_{1}(SO(3))=\mathbb Z_{2}$，对偶群 $SU(2)$ 的中心是 $\mathbb Z_{2}$——与 (iii) 一致。$A_{1}$ 型是自对偶型：$Sp_{2}=SL_{2}$ 的对偶是 $SO_{3}=PGL_{2}$，而 $SL_{2}$ 与 $PGL_{2}$ 的根数据互为对偶。
\end{example}

\begin{remark}[什么被换掉，什么没有被换掉]
对偶只看根数据，与 $G$ 的 $\mathbb C$-结构或紧实型的选取无关；$G$ 与 ${}^{L}G$ 的 Weyl 群相同（$\Delta$ 的 Weyl 群与 $\Delta^{\vee}$ 的 Weyl 群是同一个抽象群，只是作用在互为对偶的格上），维数也相同，即 $\dim G=\dim{}^{L}G$，这从生成元次数在对偶下不变看出。但整体结构互换：$G$ 单连通时 ${}^{L}G$ 是伴随型，反之亦然，故两者一般不同构；$E_{8}$、$GL_{n}$ 与 $A_{1}$ 型是自对偶的例外。又 $G$ 的中心与 ${}^{L}G$ 的基本群是同一个有限阿贝尔群的两面，把「离散通量」与「中心的 $\theta$ 角」对起来，正是 §7 的出发点。
\end{remark}

\statusline{数学陈述是定理：根数据的对换（Chevalley 的根数据分类；Langlands 1970 引入对偶群；表述见 Frenkel 2009 §1），中心与基本群的对偶是根格商的初等计算。物理陈述「${}^{L}G$ 是 $S$-对偶后的规范群」是猜想（Montonen--Olive 1977；Frenkel 2009 §4 式 (4.1)），本节只把它用作「换群」的词典条目。几何 Langlands 的猜想形式见 Frenkel 2009 §3 式 (3.1)；$GL_{n}$ 函数域情形由 Drinfeld（$n=2$）与 L.\ Lafforgue（$n>2$）证明，属另一条线索，不在本条展开。}

\dictionaryrow{它解释电磁对偶中「群被换掉」的那一半（\ref{core:em-duality}）；$n_{\mathfrak g}$ 与 $\Psi$ 的变换用它（\ref{core:gl-twist}）；几何 Langlands 对应是整本词典的终点（\ref{core:gl}）。}
\end{centry}

% ------------------------------------------------------------
% 3. 拓扑扭化与典范参数
% ------------------------------------------------------------
\begin{centry}{拓扑扭化与典范参数}{Topological Twisting and the Canonical Parameter}{gl-twist}
\oneline{把四维 $\mathcal N=4$ 超对称杨–米尔斯理论扭化成拓扑场论，得到的理论并不以扭转参数 $t\in\mathbb{CP}^{1}$ 为参数，而是以典范参数 $\Psi=\Psi(t,\tau)$ 为参数；同调意义下的理论只依赖 $\Psi$。}

\begin{definition}[扭化与扭化参数]
设 $M$ 为四维流形，$G$ 紧单。$\mathcal N=4$ 理论的 $R$-对称群为 $Spin(6)\cong SU(4)_{R}$，其超荷在 $\mathbf 4$ 中。拓扑扭化指把时空旋转群 $Spin(4)$ 换成对角子群
\[
Spin'(4)=(1\times\varkappa)\bigl(Spin(4)\bigr)\subset Spin(4)\times Spin(6),\qquad
\varkappa:Spin(4)\to Spin(6)
\]
由 $\mathrm{diag}\bigl(SU(2)_{\ell},SU(2)_{r}\bigr)\subset SU(4)_{R}$ 给出。取左、右手不变旋量 $\varepsilon_{\ell}$、$\varepsilon_{r}$，令
\[
\varepsilon=u\,\varepsilon_{\ell}+v\,\varepsilon_{r},\qquad Q=u\,Q_{\ell}+v\,Q_{r}\quad (u,v)\neq(0,0),
\]
并把 $t=v/u\in\mathbb{CP}^{1}\setminus\{0,\infty\}$ 称为扭转参数。要求 $Q^{2}=0$，这有两条件：(i) $\varepsilon$ 在 $Spin'(4)$ 下不变（这正是「重定自旋」的内容）；(ii) $\varepsilon^{T}\Gamma^{\mu}\varepsilon=0$ 对一切 $\mu$，即 $\{Q,Q\}=\varepsilon^{T}\Gamma^{\mu}\varepsilon\,P_{\mu}$ 中的平移项消失。两者一致：$\varepsilon^{T}\Gamma^{\mu}\varepsilon$ 是 $Spin'(4)$-不变标量，而扭化后 $\Gamma^{\mu}P_{\mu}$ 无 $Spin'(4)$-不变分量，故 (i) 推出 (ii)。最后定义典范参数
\[
\Psi=\Psi(t,\tau)=\mathrm{Re}\,\tau+i\,\mathrm{Im}\,\tau\cdot\frac{t-t^{-1}}{t+t^{-1}}
=\frac{\tau+\bar\tau}{2}+\frac{\tau-\bar\tau}{2}\cdot\frac{t-t^{-1}}{t+t^{-1}} .
\]
\end{definition}

\begin{theorem}[扭化理论只依赖 $\Psi$]\label{thm:psi}
设 $M$ 为闭四维流形，$G$ 紧单，固定 $t\in\mathbb{CP}^{1}\setminus\{0,\infty\}$ 与 $\tau\in\mathbb H$。引入辅助场（零形式 $P$ 与二形式 $H$）使超对称代数在壳外闭合，并作重标度 $\delta'_{t}=\delta_{t}/\sqrt{1+t^{2}}$。则总作用量可写成
\[
I=\delta_{t}V+\frac{i\Psi}{4\pi}\int_{M}\mathrm{Tr}\,F\wedge F ,
\]
其中 $V$ 与 $t$ 无关。因此：(i) 闭流形上无算子插入时的配分函数与 $t$ 无关；(ii) 更一般地，$\delta_{t}$-闭算符的关联函数只依赖 $(t,\tau)$ 的组合 $\Psi$。
\end{theorem}
\begin{proof}
由 $\{Q,Q\}=\Gamma^{\mu}P_{\mu}$ 与两个不变旋量算出超对称代数
$\delta_{T}^{2}\Phi=-i(u^{2}+v^{2})\mathcal L_{\sigma}(\Phi)$（KW06 §3.4 式 (3.35)--(3.36)），其中 $\mathcal L_{\sigma}$ 是 $\sigma$ 生成的规范变换。取 $(u,v)=(1,t)$ 并重标度，得 $(\delta'_{t})^{2}=-i\mathcal L_{\sigma}$，与 $t$ 无关。解 $\delta_{t}I_{0}=0$ 的部分取 $I_{0}=\delta_{\ell}\delta_{r}V$；剩余部分 $I_{1}$ 计算后（式 (3.46)）除一个拓扑项外全部与 $t$ 无关，而该拓扑项系数为 $(t-t^{-1})/(e^{2}(t+t^{-1}))$。加入式 (3.49) 的 $I_{2}$，其系数为 $-\bigl((t-t^{-1})/(e^{2}(t+t^{-1}))-i\theta/8\pi^{2}\bigr)$，恰好抵消，而该系数等于 $-\Psi/4\pi i$。于是总作用量为 $\delta_{t}V+(i\Psi/4\pi)\int\mathrm{Tr}\,F\wedge F$，(i)(ii) 随即成立。
\end{proof}


\begin{remark}[一般定义、三种扭化，以及「为什么变成拓扑场论」]
(a) {\bf 一般定义}：扭化是选一个子群 $Spin'(4)\subset Spin(4)\times Spin(6)$，它与 $Spin(4)$ 同构、在 $\mathbb R^4$ 上的作用与原 $Spin(4)$ 相同，但对 $\mathcal N=4$ 的超荷作用不同（KW06 §3.1 p.17）。$\mathcal N=4$ 恰有三种本质不同的扭化：两种与 Donaldson 理论类似——其瞬子不变量可用 Seiberg--Witten 不变量表达，其一在 Vafa--Witten 1994 中被详细研究；第三种由 Marcus 1995 等研究、在 KW06 之前没有应用，即本册用的 GL 扭化。
(b) {\bf 拓扑场论的定义}（Atiyah 1988 公理化）：配分函数与关联函数不依赖度规、只依赖光滑结构的量子场论。
(c) {\bf 从 $Q^{2}=0$ 到「拓扑」的推理}：扭化后取一个在 $Spin'(4)$ 下为标量的超荷 $Q=\varepsilon^{T}Q$，则 $\{Q,Q\}=2\varepsilon^{T}\Gamma^{\mu}\varepsilon\,P_{\mu}$ 的系数是普通数；取另一不变旋量 $\varepsilon'$ 使 $\varepsilon'^{T}\Gamma^{\mu}\varepsilon$ 可逆，则 $P_{\mu}$ 是 $\{Q,Q'\}$ 的线性组合，即平移算符 $Q$-恰当。于是 $Q$-上同调中 $P_{\mu}=0$：可观察量不感知平移与度规，只剩拓扑数据——这就是 (b) 意义下的拓扑场论。KW06 只写到「平移项消失」，本条的 (b)(c) 是按 Atiyah 定义补出的推理链，原文未如此陈述。
\end{remark}

\physside{同调意义下真正被理论依赖的量是 $\Psi$，$t$ 只是 $\mathbb{CP}^{1}$ 上的坐标。特殊点各有物理含义：$t=\pm i$ 时方程复化为平坦 $G_{\mathbb C}$ 联络的方程（$B$-模型侧，且 $t=\pm i$ 给出 $\Psi=\infty$，与 $\tau$ 无关）；$t=\pm1$ 时 $\Psi=\mathrm{Re}\,\tau$ 为实数，方程为实椭圆方程（$A$-模型侧，Kähler 形式 $\omega_{K}$ 进入）；$t=0,\infty$ 时方程退化，只剩瞬子或反瞬子。$S$-对偶在 $\Psi$ 上是分式线性变换：$T:\Psi\mapsto\Psi+1$、$S:\Psi\mapsto-1/(n_{\mathfrak g}\Psi)$（KW06 §3.5 式 (3.53)--(3.54)）；定向反转给 $\Psi\mapsto-\Psi$、$t\mapsto1/t$。KW06 §5 正是取 $G$ 侧 $t=1$、${}^{L}G$ 侧 ${}^{L}t=i$，把两侧分别读成 $A$-模型与 $B$-模型。}

\begin{example}[$\tau=i$，$t=2$]
取 $\tau=i$、$t=2$，则 $(t-t^{-1})/(t+t^{-1})=(2-\tfrac12)/(2+\tfrac12)=\tfrac{3}{5}$，故
\[
\Psi=0+i\cdot1\cdot\frac{3}{5}=\frac{3i}{5}.
\]
用反解公式 $t^{2}=-(\Psi-\bar\tau)/(\Psi-\tau)$（KW06 式 (3.52)）核对：右边 $=-(\tfrac{3i}{5}+i)/(\tfrac{3i}{5}-i)=-(\tfrac{8i}{5})/(-\tfrac{2i}{5})=4=t^{2}$。又 $t=\pm i$ 时 $t+t^{-1}=0$，故 $\Psi=\infty$，与 $\tau$ 无关；$t=\pm1$ 时 $\Psi=\mathrm{Re}\,\tau$。三个特值都可手算验证。再取 $t=\tfrac12$：$(t-t^{-1})/(t+t^{-1})=(-\tfrac32)/(\tfrac52)=-\tfrac35$，故 $\Psi=-\tfrac{3i}{5}$，确在 $|t|<1$ 所对应的下半平面。
\end{example}

\statusline{严格定理的部分：消灭定理（KW06 §3.3 p.28--31，式 (3.32)--(3.33)），它断言「一个 $t$ 的解对一切 $t$ 都解，且必来自平坦 $G_{\mathbb C}$ 联络」；扭化与 $Q^{2}=0$ 的两条件见 §3.1 p.17--18（三种扭化见 Vafa--Witten 1994）。「理论只依赖 $\Psi$」是论证（argued）：KW06 §3.4 p.31--35 构造作用量 (3.46)--(3.49)，§3.5 p.35--37 用重标度 $\delta'_{t}$ 证明多重态的变换律与 $t$ 无关；缺的一步是把 $Q$-上同调类的 $t$-依赖（原文只对闭流形无算子插入的配分函数断言「与 $t$ 无关」，p.35）严格化为对 $\Psi$ 的依赖。$S$ 在 $\Psi$ 上的形式见式 (3.53)--(3.54)，Frenkel 2009 §4 式 (4.2) 写作 $t\mapsto(\tau/|\tau|)t$（与原文差一符号，被 $t\sim-t$ 吸收）。}

\dictionaryrow{$S$-对偶作用在 $\Psi$ 上的形式给出 $n_{\mathfrak g}$ 的真正用武之地（\ref{core:em-duality}、\ref{core:lg}）；$t=1$ 与 $t=i$ 两值分别给出 Hitchin 模空间上的 $A$-模型与 $B$-模型（\ref{core:hitchin}、\ref{core:hitchin-fibration}）。}
\end{centry}

% ------------------------------------------------------------
% 4. Hitchin 模空间与超凯勒结构
% ------------------------------------------------------------
\begin{centry}{Hitchin 模空间与超凯勒结构}{Hitchin Moduli Space and Hyperkähler Structure}{hitchin}
\oneline{紧黎曼面上 Hitchin 方程的解模空间同时是 Higgs 丛模空间、平坦 $G_{\mathbb C}$-丛模空间，以及二维 sigma 模型的靶空间；超凯勒结构 $I,J,K$ 是「同一个空间、三种读法」的范本。}

\begin{definition}[Hitchin 方程与 Hitchin 模空间]
设 $C$ 为亏格 $g\ge2$ 的紧黎曼面，$G$ 为紧单李群，$E\to C$ 为主 $G$-丛，$A$ 为 $E$ 上联络，$\varphi\in\Omega^{1}(C,\mathrm{ad}E)$ 为伴随值一形式（Higgs 场）。以 $D=d_{A}$ 记协变外微分，$\ast$ 为 $C$ 上二维 Hodge 星，Hitchin 方程是
\[
F_{A}-\varphi\wedge\varphi=0,\qquad D\varphi=0,\qquad D\ast\varphi=0 .
\]
Hitchin 模空间 $\mathcal M_{H}(G,C)$ 定义为这些方程的解模去规范变换群所得的模空间。$g\ge2$ 时它是复流形，复维数为
\[
\dim_{\mathbb C}\mathcal M_{H}=(2g-2)\dim G ;
\]
其光滑点对应不可约解，奇点对应可约解（即存在连续未破缺规范对称性的解）。
\end{definition}

\begin{theorem}[Hitchin 1987a；非阿贝尔 Hodge 对应]
$\mathcal M_{H}(G,C)$ 上存在三个复结构 $I,J,K$，满足四元数关系
\[
I^{2}=J^{2}=K^{2}=IJK=-1,
\]
以及相应的 Kähler 形式 $\omega_{I},\omega_{J},\omega_{K}$ 与全纯辛形式
$\Omega_{I}=\omega_{J}+i\omega_{K}$ 及其循环置换。并且：
(i) 在复结构 $I$ 下，$\mathcal M_{H}$ 是半稳定 Higgs 丛 $(E,\varphi)$（$\varphi\in H^{0}(C,K_{C}\otimes\mathrm{ad}E)$）的模空间；
(ii) 在复结构 $J$ 下，$\mathcal M_{H}$ 是平坦 $G_{\mathbb C}$-丛的模空间 $Y(G,C)=\mathrm{Hom}(\pi_{1}(C),G_{\mathbb C})/G_{\mathbb C}$；
(iii) $K=IJ$；$J$ 与 $\Omega_{J}$ 不依赖 $C$ 的复结构，$\omega_{I}$ 与 $\omega_{K}$ 也不依赖。
\end{theorem}
\begin{proof}
在 $(A,\varphi)$ 的无穷维仿射空间 $W$ 上取平坦度规（KW06 §4.1 式 (4.2)），由余切空间上的转置变换定义 $I,J,K$（式 (4.3)--(4.5)），验证四元数关系与「度规在每个复结构下为 $(1,1)$ 型」，即得 $W$ 上的平坦超凯勒结构。由式 (4.8) 定义 $\omega_{I},\omega_{J},\omega_{K}$，式 (4.11) 写出三个 moment map $\mu_{I},\mu_{J},\mu_{K}$，逐项验证 Hitchin 方程等价于 $\mu_{I}=\mu_{J}=\mu_{K}=0$。按超凯勒商定理（Hitchin--Karlhede--Lindstr\"om--Ro\v{c}ek 1987），$\mathcal M_{H}=W////\mathcal G$ 仍是超凯勒流形。最后由 $I$-全纯线性函数 $A_{\bar z},\varphi_{z}$ 读出半稳定 Higgs 丛、由 $J$-全纯函数 $\mathcal A=A+i\varphi$ 读出平坦 $G_{\mathbb C}$ 联络，配合 Hitchin 1987、Corlette 1988、Simpson 1988 的非阿贝尔 Hodge 对应得到 (i)--(iii)。
\end{proof}

\physside{把四维扭化理论放在 $M=\Sigma\times C$ 上并令 $C$ 远小于 $\Sigma$，低能得到 $\Sigma$ 上的二维 sigma 模型，靶空间正是 $\mathcal M_{H}$（KW06 §4）。四维规范耦合给出靶空间数据：度规 $ds^{2}_{g}=(\mathrm{Im}\,\tau)ds^{2}$（式 (4.15)），$B$-场 $B=-(\mathrm{Re}\,\tau)\omega_{I}$（式 (4.19)）。BPS 方程就是式 (4.1)，而它等价于三个 moment map 同时为零，因此对三个辛形式都成立。三个复结构的用途不同：$J$ 下 $\Omega_{J}$ 是全纯辛形式、对应的二维理论是 $B$-模型；$K$ 下 $\omega_{K}$ 给出 $A$-模型的辛形式；至于 $\Omega_{K}$，它是 $K$ 下的全纯辛形式（$\Omega_{K}=\omega_{I}+i\omega_{J}$）。由于 $(4,4)$ 超对称，靶空间必为超凯勒商，这正是 $\mathcal M_{H}$ 超凯勒性的物理来源。}

\begin{example}[$G=SU(2)$，$g=2$]
$\dim_{\mathbb C}\mathcal M_{H}=(2g-2)\dim G=2\cdot3=6$。$\varphi=0$ 的子空间是半稳定 $SU(2)$-丛的模空间，复维数 $3(g-1)=3$（$g=2$ 时即 $3$）；固定丛 $E$ 后 $\varphi\in H^{0}(C,K_{C}\otimes\mathrm{ad}E)$ 的复维数同样是 $3(g-1)=3$（$g=2$、$\deg=0$ 时由 Serre 对偶算得 $h^{0}(K\otimes\mathrm{ad}E)=3$）。两边相加 $3+3=6$，与 $I$ 下的分解一致。$J$ 侧：平坦 $G_{\mathbb C}=SL(2,\mathbb C)$-联络的模空间复维数 $6g-6=6$、实维 $12$，与 $\dim_{\mathbb R}\mathcal M_{H}=12$ 相同；其中紧 $SU(2)$ 平坦联络的模空间实维数同为 $6g-6=6$，是 $\varphi$ 满足实条件的实子流形。两个复结构共用同一个 $12$ 维实流形（$I$ 下的复维 $6$ 与 $J$ 下的复维 $6$ 在同一实流形上）。
\end{example}

\statusline{定理：超凯勒性（Hitchin 1987a）；超凯勒商构造（Hitchin--Karlhede--Lindstr\"om--Ro\v{c}ek 1987）；非阿贝尔 Hodge 对应（Hitchin 1987、Corlette 1988、Simpson 1988）。维数公式 $\dim_{\mathbb C}\mathcal M_{H}=(2g-2)\dim G$ 见 KW06 §4 引言 p.38--39、式 (4.1) 在 p.39。物理推导：$\mathcal M_{H}$ 是紧化到 $C$ 后二维 sigma 模型的靶空间，属 KW06 §4 的物理论证（p.37--48），其严格性依赖 $(4,4)$ 超对称与超凯勒商定理。}

\dictionaryrow{它给出 Hitchin 纤维化的总空间（\ref{core:hitchin-fibration}）；它的三个复结构解释扭转参数 $t=1$ 与 $t=i$ 为何分别对应 $A$-模型与 $B$-模型（\ref{core:gl-twist}）；对偶群的 Hitchin 模空间是镜像对称的另一侧（\ref{core:lg}）。}
\end{centry}

% ------------------------------------------------------------
% 5. Hitchin 纤维化与可积系统
% ------------------------------------------------------------
\begin{centry}{Hitchin 纤维化与可积系统}{Hitchin Fibration and Integrable Systems}{hitchin-fibration}
\oneline{$\mathcal M_{H}$ 到「Casimir 谱」的映射把超凯勒流形对半切开：底空间是特征多项式的空间，一般纤维是复环面；这就是 Hitchin 1987 的完全可积系统，也是「$S$-对偶归结为 $T$-对偶」的舞台。}

\begin{definition}[Hitchin 映射]
设 $G$ 紧半单，$\mathfrak g=\mathrm{Lie}(G)$，$r=\operatorname{rank}G$。取 $(\mathrm{Fun}\,\mathfrak g)^{G}$ 的一组齐次生成元 $P_{1},\dots,P_{r}$，次数为 $d_{1},\dots,d_{r}$，它们满足
\[
\sum_{i=1}^{r}(2d_{i}-1)=\dim G .
\]
Hitchin 映射定义为
\[
\pi:\ \mathcal M_{H}(G,C)\longrightarrow \mathcal B:=\bigoplus_{i=1}^{r}H^{0}\bigl(C,K_{C}^{\otimes d_{i}}\bigr),
\qquad
(E,\varphi)\longmapsto\bigl(P_{1}(\varphi),\dots,P_{r}(\varphi)\bigr).
\]
更典范地，$\mathcal B=H^{0}\bigl(C,(\mathfrak g/\!/G)_{K_{C}}\bigr)$，其中 $(\mathfrak g/\!/G)_{K_{C}}=K_{C}\times_{\mathbb C^{\times}}\mathfrak g/\!/G$ 是由缩放 $\varphi\mapsto\lambda\varphi$ 扭曲的向量丛。$\mathcal B$ 称为 Hitchin 基。
\end{definition}

\begin{theorem}[Hitchin 1987b：完全可积性]
设 $G$ 紧半单、$C$ 亏格 $g\ge2$。则：(i) $\pi$ 在复结构 $I$ 下全纯且满射；(ii) 一般纤维是复环面（紧复阿贝尔簇，或它的主齐性空间），复维数等于 $\dim_{\mathbb C}\mathcal B$；(iii) 维数公式
\[
\dim_{\mathbb C}\mathcal B=\sum_{i=1}^{r}(2d_{i}-1)(g-1)=(g-1)\dim G=\frac12\dim_{\mathbb C}\mathcal M_{H};
\]
(iv) 纤维是 $\Omega_{I}$ 的 Lagrange 子流形（对 $\omega_{J},\omega_{K}$ 亦然）。因此 $\pi$ 是代数完全可积系统（Lagrange 纤维化）：$\mathcal M_{H}$ 上有 $\frac12\dim_{\mathbb C}\mathcal M_{H}$ 个对 Poisson 括号两两交换的全纯函数。
\end{theorem}
\begin{proof}
由 Riemann--Roch，$\dim H^{0}(C,K_{C}^{\otimes n})=(2n-1)(g-1)$（$n\ge1$），配合 $\sum(2d_{i}-1)=\dim G$ 即得 (iii)。可积性：在 $W$ 上 $\varphi_{z}$ 与 $A_{\bar z}$ 互为辛共轭，故 $\mathrm{SU}(2)$ 情形的 Hamilton 量 $H_{a}=\int_{C}\alpha_{a}\wedge\mathrm{Tr}\,\varphi^{2}$（$\alpha_{a}$ 取遍 $H^{1}(C,T_{C})$ 的基，KW06 式 (4.26)）只含 $\varphi_{z}$，彼此 Poisson 交换；辛约化保持 Poisson 交换性，于是 $\mathcal B$ 上的线性函数给出所需的交换函数族。Hitchin 1987b 进一步证明：一般纤维是椭圆方程组的解，故紧；(iv) 的 Lagrange 性来自 $\frac12\dim$ 个独立交换函数的公共水平集。(ii) 的环面性不能只由「紧 + Lagrange」得到（一般 K\"ahler 流形上的紧 Lagrange 子流形不必是环面），Hitchin 用谱覆盖与 Prym 构造把它显式证成复阿贝尔簇的主齐性空间；KW06 §4.3 明言本文不用该构造，故原文对环面性的论证是间接的：它只用纤维紧性，未使用谱覆盖构造。
\end{proof}

\physside{纤维的物理含义：在 $\mathcal B$ 的一般点处 $G$ 破缺到极大环面 $U(1)^{r}$，纤维参数化自由 $U(1)^{r}$ 规范场（平坦 $U(1)^{r}$ 联络，或等价地谱曲线的 Jacobi/Prym 簇）；底空间坐标是 Casimir 不变量的期望值 $\langle P_{i}(\varphi)\rangle$。$S$-对偶在这里变成镜像对称：$G$ 与 ${}^{L}G$ 的两个 Hitchin 纤维化共享同一个底空间（Chevalley 定理给 $\mathfrak g/\!/G=\mathfrak h/\!/W$，${}^{L}\mathfrak h=\mathfrak h^{*}$，再用不变双线性型 $\kappa_{0}$ 等同 $\mathcal B$ 与 ${}^{L}\mathcal B$），一般纤维互为对偶阿贝尔簇，镜像对称由纤维上的 $T$-对偶实现（SYZ 图像；Harvey--Moore--Strominger 1995 称之为「S-对偶约化为 T-对偶」）。}

\begin{example}[$G=SU(2)$，$g=2$]
$\dim_{\mathbb C}\mathcal M_{H}=6$；$\mathcal B=H^{0}(C,K_{C}^{2})$，$\dim_{\mathbb C}\mathcal B=3g-3=3$，恰为一半。一般纤维是复 $3$ 维环面（实维 $6$）。用谱曲线核对：$\det(\varphi-\lambda)=0$ 在 $T^{*}C$ 中给出 $C$ 的二重覆盖 $\Sigma$，亏格
\[
g_{\Sigma}=N^{2}(g-1)+1=2^{2}\cdot1+1=5 ,
\]
而 $\mathrm{Prym}(\Sigma/C)$ 的维数为 $g_{\Sigma}-g=5-2=3$，正好等于 $3g-3=3$，也等于 $\frac12\dim_{\mathbb C}\mathcal M_{H}$。再核一个秩更高的例子：$G=SU(3)$、$g=2$ 时 $d_{1}=2,d_{2}=3$，$\dim_{\mathbb C}\mathcal B=3+5=8$，而 $\dim_{\mathbb C}\mathcal M_{H}=2\cdot8=16$，纤维是复 $8$ 维环面，同样对半。
\end{example}

\begin{remark}[易错点]
不要把一般纤维写成 $T^{2r}$（$r=\operatorname{rank}G$）：纤维的复维数是 $(g-1)\dim G$（$G=SU(2)$ 时即 $3g-3$），它只依赖 $\dim G$ 与 $g$，与 rank 无关。以 $G=SU(3)$、$g=2$ 为例，$\dim G=8$、$r=2$，正确的复维数是 $8$；而 $T^{2r}$ 的复维数只有 $r=2$，相差四倍。纤维的实维数是 $2(g-1)\dim G$，恰为 $\dim_{\mathbb R}\mathcal M_{H}=2(2g-2)\dim G$ 的一半。另一点：$\mathcal B$ 与 ${}^{L}\mathcal B$ 的等同依赖 $\kappa_{0}$ 的选择，不是绝对规范。
\end{remark}

\statusline{定理：完全可积性（Hitchin 1987b，Duke Math.\ J.\ 54）；维数公式 (4.24)--(4.25) 是初等计算（KW06 §4.3 p.53--54）。纤维是环面这一点由 Hitchin 1987b 用谱覆盖证明；原文 §4.3 p.54--55 已明确声明本文不用该构造——「在本文中，我们将在不使用谱覆盖构造的前提下尽量推进」——故本节对环面性的论证是间接的（原文未走谱覆盖路线，只能用紧性加一般性论证得到）。物理陈述「纤维 = 自由 $U(1)^{r}$ 规范场、$S$-对偶 = $T$-对偶」是物理论证（Witten 2008 §3、§5.1；Frenkel 2009 §5.1；Harvey--Moore--Strominger 1995），非定理。}

\dictionaryrow{它是 Hitchin 模空间（\ref{core:hitchin}）到「Casimir 谱」的投影，也是 $S$-对偶与镜像对称同一化的关键（\ref{core:gl-twist}）；对偶群的 Hitchin 纤维化与它共享底空间（\ref{core:lg}）。}
\end{centry}
